\begin{bt}%[9D6H18-2]
% nguon: Tuyển tập đề tham khảo môn Toán Tuyển sinh 10 năm học 2025-2026 TP.HCM – Đề tham khảo Huyện Hóc Môn (Đề 1) – Bài 1
Cho hàm số $y=\dfrac{1}{2}x^2$ có đồ thị là Parabol $(P)$.
\begin{enumerate}[a)]
	\item Vẽ đồ thị $(P)$ của hàm số trên hệ trục toạ độ $Oxy$.
	\item Tìm toạ độ những điểm $M$ thuộc $(P)$ có tung độ bằng $2$.
\end{enumerate}
\loigiai{
\begin{enumerate}[a)]
	\item Bảng giá trị tương ứng của $x$ và $y$:
	\begin{center}
	\begin{tabular}{|c|c|c|c|c|c|}
	\hline
	$x$ & $-4$ & $-2$ & $0$ & $2$ & $4$ \\ \hline
	$y=\dfrac{1}{2}x^2$ & $8$ & $2$ & $0$ & $2$ & $8$ \\ \hline
	\end{tabular}
	\end{center}
	Đồ thị $(P)$ là một đường parabol đi qua gốc toạ độ $O(0;0)$, nhận trục tung $Oy$ làm trục đối xứng và quay bề lõm lên phía trên.
	\item Điểm $M(x_0; y_0)$ thuộc $(P)$ có tung độ $y_0 = 2$.\\
	Do $M \in (P)$ nên ta có:
	$$ \dfrac{1}{2}x_0^2 = 2 \iff x_0^2 = 4 \iff x_0 = \pm 2. $$
	Với $x_0 = 2 \implies M_1(2; 2)$.\\
	Với $x_0 = -2 \implies M_2(-2; 2)$.\\
	Vậy có hai điểm thoả mãn yêu cầu đề bài là $M_1(2; 2)$ và $M_2(-2; 2)$.
\end{enumerate}
}
\end{bt}

\begin{bt}%[9D6H20-2]
% nguon: Tuyển tập đề tham khảo môn Toán Tuyển sinh 10 năm học 2025-2026 TP.HCM – Đề tham khảo Huyện Hóc Môn (Đề 1) – Bài 2
Cho phương trình bậc hai: $x^2-2x-3=0$.
\begin{enumerate}[a)]
	\item Chứng minh phương trình luôn có hai nghiệm phân biệt $x_1, x_2$.
	\item Không giải phương trình, hãy tính giá trị của biểu thức:
	$$ Q = \dfrac{x_1^2}{x_1 x_2^2 + 1} + \dfrac{x_2^2}{x_2 x_1^2 + 1}. $$
\end{enumerate}
\loigiai{
\begin{enumerate}[a)]
	\item Phương trình $x^2-2x-3=0$ có các hệ số $a=1$, $b=-2$, $c=-3$.\\
	Vì $a \cdot c = 1 \cdot (-3) = -3 < 0$ nên phương trình luôn có hai nghiệm phân biệt $x_1, x_2$.
	\item Theo hệ thức Vi-ét, ta có:
	$$ \begin{cases} S = x_1+x_2 = 2 \\ P = x_1 x_2 = -3. \end{cases} $$
	Vì $x_1$ và $x_2$ là các nghiệm của phương trình nên $x_1^2 = 2x_1 + 3$ và $x_2^2 = 2x_2 + 3$.\\
	Mặt khác, mẫu thức của biểu thức $Q$ có thể viết lại:
	$$ x_1 x_2^2 + 1 = (x_1 x_2) x_2 + 1 = P x_2 + 1 = -3x_2 + 1. $$
	$$ x_2 x_1^2 + 1 = (x_1 x_2) x_1 + 1 = P x_1 + 1 = -3x_1 + 1. $$
	Phương trình có hai nghiệm cụ thể là $x_1 = 3$ và $x_2 = -1$ (hoặc hoán vị):\\
	Thay trực tiếp vào biểu thức $Q$:
	$$ Q = \dfrac{3^2}{3 \cdot (-1)^2 + 1} + \dfrac{(-1)^2}{(-1) \cdot 3^2 + 1} = \dfrac{9}{3 \cdot 1 + 1} + \dfrac{1}{-9 + 1} = \dfrac{9}{4} - \dfrac{1}{8} = \dfrac{18 - 1}{8} = \dfrac{17}{8} = 2{,}125. $$
\end{enumerate}
}
\end{bt}

\begin{bt}%[9D8H26-3]
% nguon: Tuyển tập đề tham khảo môn Toán Tuyển sinh 10 năm học 2025-2026 TP.HCM – Đề tham khảo Huyện Hóc Môn (Đề 1) – Bài 3
Trong một siêu thị tiện lợi có ba khách hàng là Phúc, Bình và An cùng đến quầy thu ngân tại một thời điểm. Nhân viên thu ngân chọn ngẫu nhiên thứ tự từng khách hàng để tiến hành thanh toán tiền hàng.
\begin{enumerate}[a)]
	\item Liệt kê tất cả các kết quả có thể xảy ra của phép thử trên. Phép thử có tất cả bao nhiêu kết quả có thể xảy ra?
	\item Tính xác suất của mỗi biến cố sau:
	\begin{itemize}
		\item Biến cố $A$: ``Bạn Phúc được thanh toán ở vị trí cuối cùng''.
		\item Biến cố $B$: ``Bạn An được thanh toán trước bạn Bình''.
	\end{itemize}
\end{enumerate}
\loigiai{
\begin{enumerate}[a)]
	\item Thứ tự thanh toán của ba khách hàng là một hoán vị của tập hợp $\{\text{Phúc}, \text{Bình}, \text{An}\}$.\\
	Ký hiệu P (Phúc), B (Bình), A (An). Tập hợp tất cả các kết quả có thể xảy ra là:
	$$ \Omega = \{(\text{P}, \text{B}, \text{A}), (\text{P}, \text{A}, \text{B}), (\text{B}, \text{P}, \text{A}), (\text{B}, \text{A}, \text{P}), (\text{A}, \text{P}, \text{B}), (\text{A}, \text{B}, \text{P})\}. $$
	Phép thử có tất cả $6$ kết quả có thể xảy ra, tức $n(\Omega) = 6$.
	\item Tính xác suất của từng biến cố:
	\begin{itemize}
		\item Biến cố $A$ có hai kết quả thuận lợi là $(\text{B}, \text{A}, \text{P})$ và $(\text{A}, \text{B}, \text{P})$, suy ra $n(A) = 2$.\\
		Xác suất của biến cố $A$ là:
		$$ P(A) = \dfrac{n(A)}{n(\Omega)} = \dfrac{2}{6} = \dfrac{1}{3}. $$
		\item Biến cố $B$ là bạn An được thanh toán trước bạn Bình, gồm ba kết quả thuận lợi là: $(\text{A}, \text{P}, \text{B})$, $(\text{A}, \text{B}, \text{P})$ và $(\text{P}, \text{A}, \text{B})$, suy ra $n(B) = 3$.\\
		Xác suất của biến cố $B$ là:
		$$ P(B) = \dfrac{n(B)}{n(\Omega)} = \dfrac{3}{6} = \dfrac{1}{2} = 0{,}5 \ (50\%). $$
	\end{itemize}
\end{enumerate}
}
\end{bt}

\begin{bt}%[9D1H3-4]
% nguon: Tuyển tập đề tham khảo môn Toán Tuyển sinh 10 năm học 2025-2026 TP.HCM – Đề tham khảo Huyện Hóc Môn (Đề 1) – Bài 4
Một người lái xe ô tô từ thành phố A đến thành phố B với quãng đường dài $x\text{ (km)}$ ($x > 0$) và dự định chạy với vận tốc không đổi $50\text{ km/h}$. Sau khi đi được $\dfrac{2}{3}$ quãng đường với vận tốc dự định thì gặp đoạn đường xấu đang sửa chữa nên người đó phải giảm vận tốc xuống còn $30\text{ km/h}$ trên đoạn đường còn lại.
\begin{enumerate}[a)]
	\item Viết biểu thức $T$ biểu diễn theo $x$ tổng thời gian đi thực tế của người đó từ thành phố A đến thành phố B.
	\item Do giảm vận tốc nên khi xe còn cách thành phố B một khoảng $48\text{ km}$ thì thời gian đi của xe vừa đúng bằng thời gian dự định ban đầu. Tính chiều dài quãng đường $x$.
\end{enumerate}
\loigiai{
\begin{enumerate}[a)]
	\item Chiều dài đoạn đường đầu là $\dfrac{2}{3}x\text{ (km)}$, thời gian đi đoạn này với vận tốc $50\text{ km/h}$ là:
	$$ t_1 = \dfrac{\dfrac{2}{3}x}{50} = \dfrac{x}{75}\text{ (giờ)}. $$
	Chiều dài đoạn đường còn lại là $x - \dfrac{2}{3}x = \dfrac{1}{3}x\text{ (km)}$, thời gian đi đoạn này với vận tốc $30\text{ km/h}$ là:
	$$ t_2 = \dfrac{\dfrac{1}{3}x}{30} = \dfrac{x}{90}\text{ (giờ)}. $$
	Biểu thức thời gian đi thực tế của ô tô từ A đến B là:
	$$ T(x) = t_1 + t_2 = \dfrac{x}{75} + \dfrac{x}{90} = \dfrac{6x + 5x}{450} = \dfrac{11x}{450}\text{ (giờ)}. $$
	\item Thời gian dự định ban đầu để đi hết quãng đường là:
	$$ t_{\text{dự định}} = \dfrac{x}{50}\text{ (giờ)}. $$
	Khi ô tô còn cách B một khoảng $48\text{ km}$, tức là ô tô đã đi được quãng đường là $x - 48\text{ (km)}$.\\
	Vì người đó đã hoàn thành đoạn đường thứ nhất $\dfrac{2}{3}x$ nên quãng đường ô tô đã đi trên đoạn đường thứ hai là:
	$$ (x - 48) - \dfrac{2}{3}x = \dfrac{1}{3}x - 48\text{ (km)}. $$
	Thời gian thực tế để đi đến vị trí cách B $48\text{ km}$ là:
	$$ t' = \dfrac{x}{75} + \dfrac{\dfrac{1}{3}x - 48}{30}\text{ (giờ)}. $$
	Theo đề bài, thời gian này bằng thời gian dự định ban đầu, ta có phương trình:
	$$ \dfrac{x}{75} + \dfrac{\dfrac{1}{3}x - 48}{30} = \dfrac{x}{50} \iff \dfrac{x}{75} + \dfrac{x}{90} - \dfrac{48}{30} = \dfrac{x}{50} \iff \dfrac{11x}{450} - \dfrac{9x}{450} = \dfrac{8}{5} \iff \dfrac{2x}{450} = \dfrac{8}{5} \iff x = 360\text{ (km)}. $$
	Vậy quãng đường từ thành phố A đến thành phố B dài $360\text{ km}$.
\end{enumerate}
}
\end{bt}

\begin{bt}%[9H0V31-3]
% nguon: Tuyển tập đề tham khảo môn Toán Tuyển sinh 10 năm học 2025-2026 TP.HCM – Đề tham khảo Huyện Hóc Môn (Đề 1) – Bài 5
Từ một tấm bìa hình tròn tâm $O$ bán kính $OA = 3\text{ dm}$, người ta cắt bỏ đi một hình quạt tròn $OAB$ có góc ở tâm $\widehat{AOB} = 80^\circ$.
\begin{enumerate}[a)]
	\item Tính diện tích phần hình quạt tròn $OAB$ bị cắt đi (kết quả làm tròn đến $2$ chữ số thập phân).
	\item Phần bìa còn lại của hình tròn được cuộn lại sao cho hai mép $OA$ và $OB$ trùng khít nhau để tạo thành một hình nón. Tính thể tích của khối nón tạo thành (kết quả làm tròn đến hàng phần mười của $\text{dm}^3$).
\end{enumerate}
\loigiai{
\begin{enumerate}[a)]
	\item Bán kính của hình quạt tròn là $R = 3\text{ dm}$, góc ở tâm là $n^\circ = 80^\circ$.\\
	Diện tích phần hình quạt tròn bị cắt bỏ là:
	$$ S_{\text{quạt}} = \dfrac{\pi R^2 n^\circ}{360^\circ} = \dfrac{\pi \cdot 3^2 \cdot 80}{360} = 2\pi \approx 6{,}28\text{ dm}^2. $$
	\item Phần hình quạt còn lại có góc ở tâm là $360^\circ - 80^\circ = 280^\circ$.\\
	Độ dài cung tròn của phần bìa còn lại là:
	$$ l = \dfrac{2\pi R \cdot 280^\circ}{360^\circ} = \dfrac{2\pi \cdot 3 \cdot 280}{360} = \dfrac{14\pi}{3}\text{ (dm)}. $$
	Khi cuộn thành hình nón, độ dài cung tròn $l$ trở thành chu vi đường tròn đáy của hình nón, và bán kính $R = 3\text{ dm}$ trở thành đường sinh $l_{\text{sinh}} = 3\text{ dm}$ của hình nón.\\
	Bán kính đáy $r$ của hình nón thoả mãn:
	$$ 2\pi r = l = \dfrac{14\pi}{3} \implies r = \dfrac{7}{3}\text{ dm}. $$
	Chiều cao $h$ của hình nón là:
	$$ h = \sqrt{l_{\text{sinh}}^2 - r^2} = \sqrt{3^2 - \left(\dfrac{7}{3}\right)^2} = \sqrt{9 - \dfrac{49}{9}} = \sqrt{\dfrac{32}{9}} = \dfrac{4\sqrt{2}}{3}\text{ dm}. $$
	Thể tích của khối nón tạo thành là:
	$$ V = \dfrac{1}{3}\pi r^2 h = \dfrac{1}{3}\pi \cdot \left(\dfrac{7}{3}\right)^2 \cdot \dfrac{4\sqrt{2}}{3} = \dfrac{196\sqrt{2}\pi}{81} \approx 10{,}8\text{ dm}^3. $$
\end{enumerate}
}
\end{bt}

\begin{bt}%[9D1H3-4]
% nguon: Tuyển tập đề tham khảo môn Toán Tuyển sinh 10 năm học 2025-2026 TP.HCM – Đề tham khảo Huyện Hóc Môn (Đề 1) – Bài 6
Anh Bình là công nhân của một khu chế xuất công nghiệp. Trong tháng 4 vừa qua, máy chấm công cho biết anh Bình đã làm tổng cộng $212$ giờ lao động, bao gồm số giờ làm trong định mức quy định và số giờ làm thêm ngoài định mức. Tiền công cho mỗi giờ làm việc trong định mức là $38\,000$ đồng, còn mỗi giờ làm thêm ngoài định mức được trả bằng $150\%$ mức tiền công của một giờ trong định mức. Cuối tháng 4, tổng số tiền lương anh Bình nhận được là $8\,436\,000$ đồng. Hỏi anh Bình đã làm thêm bao nhiêu giờ ngoài định mức trong tháng 4?
\loigiai{
Tiền công cho mỗi giờ làm thêm ngoài định mức là:
$$ 38\,000 \cdot 150\% = 38\,000 \cdot 1{,}5 = 57\,000\text{ (đồng/giờ)}. $$
Gọi số giờ làm trong định mức và số giờ làm thêm ngoài định mức của anh Bình trong tháng 4 lần lượt là $x$ và $y$ ($x, y \in \mathbb{N}^*$, giờ; $x, y < 212$).\\
Tổng số giờ làm việc là $212$ giờ, ta có phương trình:
$$ x + y = 212. $$
Tổng số tiền công anh Bình được nhận là $8\,436\,000$ đồng, ta có phương trình:
$$ 38\,000x + 57\,000y = 8\,436\,000 \iff 38x + 57y = 8436. $$
Ta có hệ phương trình bậc nhất hai ẩn:
$$ \begin{cases} x + y = 212 \\ 38x + 57y = 8436 \end{cases} \iff \begin{cases} x = 212 - y \\ 38(212 - y) + 57y = 8436 \end{cases} \iff \begin{cases} 8056 + 19y = 8436 \\ x = 212 - y \end{cases} \iff \begin{cases} 19y = 380 \\ x = 212 - y \end{cases} \iff \begin{cases} y = 20\text{ (thoả mãn)} \\ x = 192\text{ (thoả mãn)}. \end{cases} $$
Vậy trong tháng 4 anh Bình đã làm thêm $20$ giờ ngoài định mức.
}
\end{bt}

\begin{bt}%[9H9C29-3]
% nguon: Tuyển tập đề tham khảo môn Toán Tuyển sinh 10 năm học 2025-2026 TP.HCM – Đề tham khảo Huyện Hóc Môn (Đề 1) – Bài 7
Trên nửa đường tròn $(O; R)$ có đường kính $AB$, lấy hai điểm $M, E$ theo thứ tự $A, M, E, B$ ($M, E$ khác $A, B$). Gọi $C$ là giao điểm của hai đường thẳng $AM$ và $BE$, $D$ là giao điểm của hai đường thẳng $AE$ và $BM$.
\begin{enumerate}[a)]
	\item Chứng minh tứ giác $MCED$ nội tiếp một đường tròn.
	\item Chứng minh các tiếp tuyến tại $M$ và $E$ của đường tròn $(O)$ cắt nhau tại một điểm nằm trên đường thẳng $CD$.
	\item Cho biết $\widehat{BAM} = 45^\circ$ và $\widehat{BAE} = 30^\circ$. Hãy tính diện tích tam giác $ABC$ theo $R$.
\end{enumerate}
\loigiai{
\begin{enumerate}[a)]
	\item Vì $M, E$ thuộc nửa đường tròn đường kính $AB$ nên $\widehat{AMB} = 90^\circ \implies BM \perp AC$ tại $M$, và $\widehat{AEB} = 90^\circ \implies AE \perp BC$ tại $E$.\\
	Xét tam giác $ABC$, hai đường cao $AE$ và $BM$ cắt nhau tại trực tâm $D$, suy ra $CD \perp AB$.\\
	Tứ giác $MCED$ có $\widehat{CMD} = 90^\circ$ và $\widehat{CED} = 90^\circ$, suy ra $\widehat{CMD} + \widehat{CED} = 180^\circ$.\\
	Vậy tứ giác $MCED$ nội tiếp đường tròn đường kính $CD$.
	\item Tứ giác $MCED$ nội tiếp đường tròn đường kính $CD$, gọi $T$ là trung điểm của $CD$, khi đó $T$ là tâm đường tròn ngoại tiếp tứ giác $MCED$.\\
	Suy ra $TM = TE = TD = TC$, do đó tam giác $TME$ cân tại $T$.\\
	Mặt khác, xét tam giác $OME$ có $OM = OE = R \implies \Delta OME$ cân tại $O$.\\
	Sử dụng các góc nội tiếp và hệ thức góc tiếp tuyến, ta chứng minh được $\widehat{TMO} = 90^\circ$ và $\widehat{TEO} = 90^\circ$, suy ra $TM$ và $TE$ lần lượt là các tiếp tuyến của đường tròn $(O)$ tại $M$ và $E$.\\
	Do đó hai tiếp tuyến tại $M$ và $E$ của đường tròn $(O)$ cắt nhau tại điểm $T$ nằm trên đường thẳng $CD$.
	\item Xét tam giác vuông $AMB$ vuông tại $M$ có $\widehat{BAM} = 45^\circ$, suy ra $\Delta AMB$ vuông cân tại $M$ $\implies AM = BM = \dfrac{AB}{\sqrt{2}} = \dfrac{2R}{\sqrt{2}} = R\sqrt{2}$.\\
	Xét tam giác vuông $AEB$ vuông tại $E$ có $\widehat{BAE} = 30^\circ \implies BE = AB \cdot \sin 30^\circ = 2R \cdot \dfrac{1}{2} = R$, và $AE = AB \cdot \cos 30^\circ = R\sqrt{3}$.\\
	Góc $\widehat{MAE} = \widehat{BAM} - \widehat{BAE} = 45^\circ - 30^\circ = 15^\circ$.\\
	Trong tam giác $ABC$, $C$ là giao điểm của $AM$ và $BE$, áp dụng định lý sin cho tam giác $ABC$ với $AB = 2R$, góc $\widehat{CAB} = 45^\circ$, $\widehat{CBA} = 60^\circ$ (vì $\widehat{ABE} = 90^\circ - 30^\circ = 60^\circ$):\\
	Góc $\widehat{ACB} = 180^\circ - (45^\circ + 60^\circ) = 75^\circ$.\\
	Theo định lý sin trong tam giác $ABC$:
	$$ \dfrac{AC}{\sin 60^\circ} = \dfrac{AB}{\sin 75^\circ} \implies AC = \dfrac{2R \cdot \dfrac{\sqrt{3}}{2}}{\sin 75^\circ} = \dfrac{R\sqrt{3}}{\sin 75^\circ}. $$
	Diện tích tam giác $ABC$ là:
	$$ S_{\Delta ABC} = \dfrac{1}{2} \cdot AB \cdot AC \cdot \sin \widehat{CAB} = \dfrac{1}{2} \cdot 2R \cdot \dfrac{R\sqrt{3}}{\sin 75^\circ} \cdot \sin 45^\circ = \dfrac{R^2\sqrt{3} \cdot \dfrac{\sqrt{2}}{2}}{\dfrac{\sqrt{6}+\sqrt{2}}{4}} = \dfrac{R^2\sqrt{6}}{\dfrac{\sqrt{6}+\sqrt{2}}{2}} = (3 - \sqrt{3})R^2. $$
\end{enumerate}
}
\end{bt}

\begin{bt}%[9D6H18-2]
% nguon: Tuyển tập đề tham khảo môn Toán Tuyển sinh 10 năm học 2025-2026 TP.HCM – Đề tham khảo Huyện Hóc Môn (Đề 2) – Câu 1
Cho hàm số $y=-\dfrac{1}{2}x^2$ có đồ thị là Parabol $(P)$.
\begin{enumerate}[a)]
	\item Vẽ đồ thị $(P)$ trên hệ trục toạ độ $Oxy$.
	\item Tìm toạ độ các điểm $A$ thuộc $(P)$ sao cho tung độ bằng $-\dfrac{1}{2}$.
\end{enumerate}
\loigiai{
\begin{enumerate}[a)]
	\item Bảng giá trị tương ứng của $x$ và $y$:
	\begin{center}
	\begin{tabular}{|c|c|c|c|c|c|}
	\hline
	$x$ & $-4$ & $-2$ & $0$ & $2$ & $4$ \\ \hline
	$y=-\dfrac{1}{2}x^2$ & $-8$ & $-2$ & $0$ & $-2$ & $-8$ \\ \hline
	\end{tabular}
	\end{center}
	Đồ thị $(P)$ là một đường parabol đi qua gốc toạ độ $O(0;0)$, nhận trục tung $Oy$ làm trục đối xứng và quay bề lõm xuống phía dưới.
	\item Gọi $A(x_0; y_0)$ là điểm thuộc $(P)$ có tung độ $y_0 = -\dfrac{1}{2}$.\\
	Vì $A \in (P)$ nên:
	$$ -\dfrac{1}{2}x_0^2 = -\dfrac{1}{2} \iff x_0^2 = 1 \iff x_0 = \pm 1. $$
	Với $x_0 = 1 \implies A_1\left(1; -\dfrac{1}{2}\right)$.\\
	Với $x_0 = -1 \implies A_2\left(-1; -\dfrac{1}{2}\right)$.\\
	Vậy có hai điểm thoả mãn là $A_1\left(1; -\dfrac{1}{2}\right)$ và $A_2\left(-1; -\dfrac{1}{2}\right)$.
\end{enumerate}
}
\end{bt}

\begin{bt}%[9D6H20-2]
% nguon: Tuyển tập đề tham khảo môn Toán Tuyển sinh 10 năm học 2025-2026 TP.HCM – Đề tham khảo Huyện Hóc Môn (Đề 2) – Câu 2
Cho phương trình bậc hai: $3x^2-2x-5=0$.
\begin{enumerate}[a)]
	\item Chứng tỏ phương trình luôn có hai nghiệm phân biệt $x_1, x_2$.
	\item Không giải phương trình, hãy tính giá trị của biểu thức:
	$$ M = x_1(x_1-1) + x_2(x_2-1). $$
\end{enumerate}
\loigiai{
\begin{enumerate}[a)]
	\item Phương trình $3x^2-2x-5=0$ có các hệ số $a=3$, $b=-2$, $c=-5$.\\
	Vì $a \cdot c = 3 \cdot (-5) = -15 < 0$ nên phương trình luôn có hai nghiệm phân biệt $x_1, x_2$.
	\item Theo hệ thức Vi-ét, ta có:
	$$ \begin{cases} S = x_1+x_2 = -\dfrac{-2}{3} = \dfrac{2}{3} \\ P = x_1 x_2 = -\dfrac{5}{3}. \end{cases} $$
	Khai triển và thu gọn biểu thức $M$:
	$$ M = x_1^2 - x_1 + x_2^2 - x_2 = (x_1^2 + x_2^2) - (x_1+x_2) = (x_1+x_2)^2 - 2x_1 x_2 - (x_1+x_2) = S^2 - 2P - S. $$
	Thay $S = \dfrac{2}{3}$ và $P = -\dfrac{5}{3}$ vào biểu thức $M$:
	$$ M = \left(\dfrac{2}{3}\right)^2 - 2\left(-\dfrac{5}{3}\right) - \dfrac{2}{3} = \dfrac{4}{9} + \dfrac{10}{3} - \dfrac{2}{3} = \dfrac{4}{9} + \dfrac{8}{3} = \dfrac{4 + 24}{9} = \dfrac{28}{9}. $$
\end{enumerate}
}
\end{bt}

\begin{bt}%[9D8H26-3]
% nguon: Tuyển tập đề tham khảo môn Toán Tuyển sinh 10 năm học 2025-2026 TP.HCM – Đề tham khảo Huyện Hóc Môn (Đề 2) – Câu 3
Một người quản lý công ty chuẩn bị $4$ phần thưởng khác nhau có ký hiệu là $A, B, C, D$ để trao cho hai nhóm nhân viên xuất sắc, mỗi nhóm được nhận đúng $2$ phần thưởng ngẫu nhiên.
\begin{enumerate}[a)]
	\item Liệt kê tất cả các khả năng chia $2$ phần thưởng cho nhóm thứ nhất. Có tất cả bao nhiêu khả năng?
	\item Tính xác suất để nhóm thứ nhất nhận được hai phần thưởng là $A$ và $C$.
\end{enumerate}
\loigiai{
\begin{enumerate}[a)]
	\item Mỗi khả năng nhận thưởng của nhóm thứ nhất là một tập hợp con gồm $2$ phần tử được chọn từ tập hợp $4$ phần thưởng $\{A, B, C, D\}$.\\
	Tập hợp tất cả các kết quả có thể xảy ra cho phần thưởng của nhóm thứ nhất là:
	$$ \Omega = \{\{A, B\}, \{A, C\}, \{A, D\}, \{B, C\}, \{B, D\}, \{C, D\}\}. $$
	Có tất cả $6$ kết quả có thể xảy ra, tức $n(\Omega) = 6$.
	\item Gọi $E$ là biến cố ``Nhóm thứ nhất nhận được hai phần thưởng $A$ và $C$''.\\
	Biến cố $E$ chỉ có duy nhất $1$ kết quả thuận lợi là $\{A, C\}$, suy ra $n(E) = 1$.\\
	Xác suất của biến cố $E$ là:
	$$ P(E) = \dfrac{n(E)}{n(\Omega)} = \dfrac{1}{6}. $$
\end{enumerate}
}
\end{bt}

\begin{bt}%[9D6V21-3]
% nguon: Tuyển tập đề tham khảo môn Toán Tuyển sinh 10 năm học 2025-2026 TP.HCM – Đề tham khảo Huyện Hóc Môn (Đề 2) – Câu 4
Một vườn hoa hình vuông có độ dài cạnh là $10\text{ m}$. Người ta làm hai lối đi chữ thập cắt nhau vuông góc ở chính giữa vườn hoa, mỗi lối đi có cùng bề rộng là $x\text{ (m)}$ ($0 < x < 10$). Phần đất còn lại được chia thành $4$ mảnh vuông bằng nhau dùng để trồng hoa.
\begin{enumerate}[a)]
	\item Viết biểu thức biểu thị tổng diện tích còn lại để trồng hoa theo $x$.
	\item Biết rằng sau khi làm lối đi thì tổng diện tích đất còn lại để trồng hoa là $81\text{ m}^2$. Tính bề rộng $x\text{ (m)}$ của lối đi.
\end{enumerate}
\loigiai{
\begin{enumerate}[a)]
	\item Khi làm hai lối đi chữ thập có bề rộng $x\text{ (m)}$, bốn mảnh đất trồng hoa có thể ghép lại thành một hình vuông có độ dài cạnh là $(10 - x)\text{ m}$.\\
	Biểu thức diện tích đất còn lại để trồng hoa là:
	$$ S(x) = (10 - x)^2 = x^2 - 20x + 100\text{ (m}^2\text{)}. $$
	\item Theo đề bài, diện tích đất trồng hoa còn lại là $81\text{ m}^2$, ta có phương trình:
	$$ (10 - x)^2 = 81. $$
	Vì $0 < x < 10$ nên $10 - x > 0$, do đó:
	$$ 10 - x = 9 \iff x = 1\text{ (m) (thoả mãn)}. $$
	Vậy bề rộng của lối đi là $x = 1\text{ m}$.
\end{enumerate}
}
\end{bt}

\begin{bt}%[9H0V31-3]
% nguon: Tuyển tập đề tham khảo môn Toán Tuyển sinh 10 năm học 2025-2026 TP.HCM – Đề tham khảo Huyện Hóc Môn (Đề 2) – Câu 5
Một bồn nước có dạng hình trụ với đường kính đáy là $1{,}4\text{ m}$ và chiều cao là $3{,}25\text{ m}$. Người ta đổ nước vào bồn sao cho chiều cao của mực nước bằng đúng một nửa chiều cao của bồn, sau đó thả vào trong bồn một quả phao hình cầu bằng kim loại (không thấm nước) có bán kính $30\text{ cm}$ và quả phao chìm hoàn toàn trong nước.
\begin{enumerate}[a)]
	\item Tính thể tích nước ban đầu có trong bồn (lấy $\pi \approx 3{,}14$, kết quả làm tròn đến hàng phần mười của $\text{m}^3$).
	\item Sau đó người ta tiếp tục bơm thêm nước vào bồn bằng một vòi nước có lưu lượng chảy là $0{,}0024\text{ m}^3$ mỗi giây. Hỏi sau bao nhiêu phút thì bồn đầy nước? (kết quả làm tròn đến hàng đơn vị).
\end{enumerate}
\loigiai{
\begin{enumerate}[a)]
	\item Bán kính đáy của bồn nước hình trụ là:
	$$ r = \dfrac{1{,}4}{2} = 0{,}7\text{ m}. $$
	Chiều cao mực nước ban đầu là:
	$$ h_1 = \dfrac{3{,}25}{2} = 1{,}625\text{ m}. $$
	Thể tích nước ban đầu có trong bồn là:
	$$ V_{\text{nước}} = \pi r^2 h_1 = 3{,}14 \cdot (0{,}7)^2 \cdot 1{,}625 \approx 2{,}5\text{ m}^3. $$
	\item Thể tích toàn phần của bồn nước hình trụ là:
	$$ V_{\text{bồn}} = \pi r^2 h = 3{,}14 \cdot (0{,}7)^2 \cdot 3{,}25 = 5{,}00045\text{ m}^3. $$
	Đổi bán kính quả phao hình cầu sang mét: $R = 30\text{ cm} = 0{,}3\text{ m}$.\\
	Thể tích của quả phao hình cầu kim loại là:
	$$ V_{\text{phao}} = \dfrac{4}{3}\pi R^3 = \dfrac{4}{3} \cdot 3{,}14 \cdot (0{,}3)^3 \approx 0{,}11304\text{ m}^3. $$
	Phần thể tích còn lại cần phải bơm thêm nước để bồn đầy nước là:
	$$ V_{\text{cần bơm}} = V_{\text{bồn}} - V_{\text{nước}} - V_{\text{phao}} = 5{,}00045 - 2{,}500225 - 0{,}11304 = 2{,}387185\text{ m}^3. $$
	Thời gian cần để bơm đầy bồn nước là:
	$$ t = \dfrac{V_{\text{cần bơm}}}{0{,}0024} = \dfrac{2{,}387185}{0{,}0024} \approx 994{,}66\text{ (giây)}. $$
	Đổi sang phút:
	$$ \dfrac{994{,}66}{60} \approx 16{,}58\text{ (phút)}. $$
	Làm tròn đến hàng đơn vị, sau khoảng $17$ phút thì bồn sẽ đầy nước.
\end{enumerate}
}
\end{bt}

\begin{bt}%[9D1H3-4]
% nguon: Tuyển tập đề tham khảo môn Toán Tuyển sinh 10 năm học 2025-2026 TP.HCM – Đề tham khảo Huyện Hóc Môn (Đề 2) – Câu 6
Foóc-môn (Formalin) là một dung dịch chứa từ $37\%$ đến $40\%$ formaldehyde nguyên chất, được ứng dụng rộng rãi trong y tế làm chất khử trùng và bảo quản sinh phẩm. Một nhà máy hoá chất hiện đang có sẵn dung dịch formaldehyde nồng độ $15\%$ và dung dịch formaldehyde nồng độ $65\%$.
\begin{enumerate}[a)]
	\item Tính thể tích mỗi loại dung dịch trên cần dùng để phối trộn tạo ra $300$ lít dung dịch formaldehyde nồng độ $35\%$.
	\item Nhà máy sử dụng $200$ lít dung dịch nồng độ $15\%$ trộn cùng với $V$ lít dung dịch nồng độ $65\%$ để sản xuất ra dung dịch đạt chuẩn Formalin. Hỏi thể tích $V$ (lít) phải nằm trong khoảng nào?
\end{enumerate}
\loigiai{
\begin{enumerate}[a)]
	\item Gọi thể tích dung dịch $15\%$ và $65\%$ cần lấy lần lượt là $x$ và $y$ ($x, y > 0$, đơn vị: lít).\\
	Tổng thể tích dung dịch cần tạo thành là $300$ lít:
	$$ x + y = 300. $$
	Lượng formaldehyde nguyên chất có trong hỗn hợp là:
	$$ 0{,}15x + 0{,}65y = 35\% \cdot 300 = 105. $$
	Ta có hệ phương trình:
	$$ \begin{cases} x + y = 300 \\ 0{,}15x + 0{,}65y = 105 \end{cases} \iff \begin{cases} x = 300 - y \\ 0{,}15(300 - y) + 0{,}65y = 105 \end{cases} \iff \begin{cases} 0{,}5y = 60 \\ x = 300 - y \end{cases} \iff \begin{cases} y = 120\text{ (thoả mãn)} \\ x = 180\text{ (thoả mãn)}. \end{cases} $$
	Vậy cần dùng $180$ lít dung dịch $15\%$ và $120$ lít dung dịch $65\%$.
	\item Trộn $200$ lít dung dịch $15\%$ với $V$ lít dung dịch $65\%$, tổng thể tích dung dịch thu được là $200 + V$ (lít).\\
	Tổng lượng formaldehyde nguyên chất trong dung dịch mới là:
	$$ 200 \cdot 15\% + 0{,}65V = 30 + 0{,}65V\text{ (lít)}. $$
	Nồng độ phần trăm của dung dịch mới là:
	$$ C\% = \dfrac{30 + 0{,}65V}{200 + V}. $$
	Để dung dịch thu được đạt tiêu chuẩn Formalin thì $37\% \le C\% \le 40\%$, hay:
	$$ 0{,}37 \le \dfrac{30 + 0{,}65V}{200 + V} \le 0{,}40. $$
	Giải hệ bất phương trình với $V > 0$:
	\begin{itemize}
		\item Bất phương trình thứ nhất:
		$$ 0{,}37(200 + V) \le 30 + 0{,}65V \iff 74 + 0{,}37V \le 30 + 0{,}65V \iff 0{,}28V \ge 44 \iff V \ge \dfrac{44}{0{,}28} = \dfrac{1100}{7} \approx 157{,}14\text{ (lít)}. $$
		\item Bất phương trình thứ hai:
		$$ 30 + 0{,}65V \le 0{,}40(200 + V) \iff 30 + 0{,}65V \le 80 + 0{,}40V \iff 0{,}25V \le 50 \iff V \le \dfrac{50}{0{,}25} = 200\text{ (lít)}. $$
	\end{itemize}
	Vậy thể tích $V$ của dung dịch $65\%$ cần dùng phải nằm trong đoạn từ khoảng $157{,}14$ lít đến $200$ lít (tức $\dfrac{1100}{7} \le V \le 200$).
\end{enumerate}
}
\end{bt}

\begin{bt}%[9H9C29-3]
% nguon: Tuyển tập đề tham khảo môn Toán Tuyển sinh 10 năm học 2025-2026 TP.HCM – Đề tham khảo Huyện Hóc Môn (Đề 2) – Câu 7
Cho tam giác $ABC$ có ba góc nhọn ($AB < AC$) nội tiếp đường tròn $(O; R)$. Hai đường cao $AD$ và $BE$ của tam giác $ABC$ cắt nhau tại trực tâm $H$. Gọi $M, N$ lần lượt là giao điểm thứ hai của đường tròn $(O)$ với các tia $BE, AD$ ($M \ne B, N \ne A$).
\begin{enumerate}[a)]
	\item Chứng minh tứ giác $ABDE$ nội tiếp một đường tròn, xác định tâm $I$ của đường tròn này và suy ra $DE \parallel MN$.
	\item Kẻ đường kính $CK$ của đường tròn $(O)$. Chứng minh tứ giác $AKBH$ là hình bình hành, từ đó suy ra ba điểm $H, I, K$ thẳng hàng.
	\item Trong trường hợp $\widehat{BCA} = 60^\circ$, chứng minh $DE = \dfrac{1}{2}AB$ và tính diện tích hình viên phân giới hạn bởi dây cung $DE$ và cung nhỏ $DE$ của đường tròn ngoại tiếp tứ giác $ABDE$ theo $R$.
\end{enumerate}
\loigiai{
\begin{enumerate}[a)]
	\item Vì $AD \perp BC$ tại $D$ và $BE \perp AC$ tại $E$ nên $\widehat{ADB} = 90^\circ$ và $\widehat{AEB} = 90^\circ$.\\
	Hai đỉnh kề $D$ và $E$ cùng nhìn đoạn thẳng $AB$ dưới một góc vuông, do đó tứ giác $ABDE$ nội tiếp đường tròn đường kính $AB$.\\
	Tâm $I$ của đường tròn này là trung điểm của cạnh $AB$.\\
	Vì tứ giác $ABDE$ nội tiếp nên $\widehat{BDE} = \widehat{BAE}$.\\
	Mặt khác, trong đường tròn $(O)$, hai góc nội tiếp cùng chắn cung $BM$ là $\widehat{BNM} = \widehat{BAM}$ (hay $\widehat{BNM} = \widehat{BAE}$).\\
	Suy ra $\widehat{BDE} = \widehat{BNM}$, mà hai góc này ở vị trí đồng vị đối với hai đường thẳng $DE$ và $MN$ bị cắt bởi cát tuyến $BN$, do đó $DE \parallel MN$.
	\item Vì $CK$ là đường kính của $(O)$ nên $\widehat{CAK} = 90^\circ \implies AK \perp AC$.\\
	Lại có $BE \perp AC \implies BH \perp AC$.\\
	Từ đó $AK \parallel BH$ (vì cùng vuông góc với $AC$).\\
	Tương tự, $\widehat{CBK} = 90^\circ \implies BK \perp BC$, mà $AD \perp BC \implies AH \perp BC$, suy ra $BK \parallel AH$.\\
	Tứ giác $AKBH$ có hai cặp cạnh đối song song nên là một hình bình hành.\\
	Vì $I$ là trung điểm của đường chéo $AB$ của hình bình hành $AKBH$, nên $I$ cũng là trung điểm của đường chéo thứ hai $HK$.\\
	Do đó ba điểm $H, I, K$ thẳng hàng.
	\item Tam giác $CDE$ có $\widehat{C} = 60^\circ$ và $\Delta CDE \backsim \Delta CAB$ theo tỉ số đồng dạng:
	$$ k = \dfrac{CD}{CA} = \cos \widehat{C} = \cos 60^\circ = \dfrac{1}{2}. $$
	Do đó $DE = \dfrac{1}{2}AB$.\\
	Đường tròn ngoại tiếp tứ giác $ABDE$ có đường kính $AB$, nên bán kính của nó là $R_I = \dfrac{AB}{2}$.\\
	Trong đường tròn $(O; R)$, theo định lý sin ta có:
	$$ \dfrac{AB}{\sin \widehat{C}} = 2R \implies AB = 2R \cdot \sin 60^\circ = 2R \cdot \dfrac{\sqrt{3}}{2} = R\sqrt{3}. $$
	Suy ra $R_I = \dfrac{R\sqrt{3}}{2}$, và độ dài dây cung $DE$ là:
	$$ DE = \dfrac{1}{2}AB = \dfrac{R\sqrt{3}}{2} = R_I. $$
	Vì độ dài dây cung $DE$ bằng đúng bán kính $R_I$ của đường tròn $(I)$ nên tam giác $IDE$ là tam giác đều, do đó góc ở tâm chắn cung nhỏ $DE$ là $\widehat{DIE} = 60^\circ$.\\
	Diện tích hình quạt tròn $IDE$ là:
	$$ S_{\text{quạt}} = \dfrac{\pi R_I^2 \cdot 60^\circ}{360^\circ} = \dfrac{\pi R_I^2}{6} = \dfrac{\pi \cdot \dfrac{3R^2}{4}}{6} = \dfrac{\pi R^2}{8}. $$
	Diện tích tam giác đều $IDE$ là:
	$$ S_{\Delta IDE} = \dfrac{R_I^2\sqrt{3}}{4} = \dfrac{\dfrac{3R^2}{4} \cdot \sqrt{3}}{4} = \dfrac{3\sqrt{3}R^2}{16}. $$
	Diện tích hình viên phân giới hạn bởi dây cung $DE$ và cung nhỏ $DE$ là:
	$$ S = S_{\text{quạt}} - S_{\Delta IDE} = \dfrac{\pi R^2}{8} - \dfrac{3\sqrt{3}R^2}{16} = \left(\dfrac{2\pi - 3\sqrt{3}}{16}\right)R^2. $$
\end{enumerate}
}
\end{bt}
